\section{项目摘要}
\subsection{项目背景}
随着电子商务的迅猛发展，我国快递业务量连年高速增长。据统计，2024 年全国快递业务量已突破 1400 亿件。在末端配送环节，社区快递代收点（如菜鸟驿站、丰巢快递柜、社区便利店代收点等）已成为连接快递员与收件人的关键枢纽。

然而，许多中小型代收点仍采用手工记账或简单 Excel 表格的方式进行快递管理。这种方式存在以下问题：

• 效率低下：手工查找快递耗时长，高峰期容易出现排队拥堵

• 容易出错：快递单号手工记录易发生抄写错误，导致取错件或丢件

• 统计困难：无法快速了解当前库存、各快递公司占比、货架占用情况

• 信息滞后：无法实时知道每个快递的状态（待取/已取/已退）

因此，开发一套轻量级、易于部署的快递代收点管理系统，对于提升社区快递末端服务效率具有现实意义。

\subsection{项目目标}
本项目旨在使用 C 语言开发一套基于控制台的快递代收点管理系统，实现以下核心目标：

1. 信息化管理：将快递入库、查询、取件、退件等日常操作电子化，替代手工记录

2. 状态追踪：实时跟踪每个快递的状态（待取/已取/已退），避免错取、漏取

3. 数据统计：提供按状态、快递公司、货架等多维度的统计功能，辅助运营决策

4. 数据持久化：所有快递信息以文本文件存储，程序重启后数据不丢失

5. 操作简便：通过菜单交互完成所有操作，无需数据库等复杂环境配置

\subsection{环境要求}
\subsection{功能要求}
本系统需实现以下核心功能模块：

1. 快递录入：录入新到达的快递信息，自动校验单号唯一性

2. 快递查询：支持按单号精确查询和按姓名模糊查询

3. 取件管理：收件人取件时标记为“已取”

4. 退件处理：超期未取或拒收的快递标记为“已退”

5. 信息修改：支持修改已录入快递的各项信息

6. 数据统计：统计各状态数量、各公司占比、各货架分布

此外还包括浏览全部快递和删除快递两个辅助功能，以及文件持久化自动保存与加载。

\subsection{关键技术}
\section{需求分析}
\subsection{数据需求}
\subsubsection{快递数据结构}
每个快递记录包含以下七个字段，在程序中定义为 Package 结构体：

\subsubsection{存储格式}
数据以空格分隔的文本行存储在 packages.txt 文件中，每行一条记录。程序启动时自动加载，退出时自动保存。最大记录数 100 条（宏 MAX\_PACKAGES 可调）。

\subsection{功能需求说明}
快递录入(add\_package)

• 输入:单号、姓名、电话、公司、货架、日期

• 处理:检查上限→唯一性→设待取→追加→保存

• 输出:录入成功+总数

快递查询(query\_package)

• 输入:选择方式(1-单号/2-姓名)

• 处理:线性查找/strstr模糊匹配

• 输出:详情或未找到

取件管理(pickup\_package)

• 输入:单号

• 处理:查找→检查状态→更新已取→保存

• 输出:取件成功

退件处理(return\_package)

• 输入:单号

• 处理:判断状态→更新已退→保存

信息修改(modify\_package)

• 输入:单号+逐字段确认

• 处理:逐字段询问,回车保留原值

• 输出:修改成功

数据统计(show\_statistics)

• 处理:按状态/公司/货架去重统计

• 输出:分类统计表格

\subsection{功能模块划分}
1. 快递录入—add\_package()

2. 快递查询—query\_package()

3. 取件管理—pickup\_package()

4. 退件处理—return\_package()

5. 信息修改—modify\_package()

6. 数据统计—show\_statistics()

7. 浏览全部—browse\_all()

8. 删除快递—delete\_package()

9. 文件读写—load\_data()/save\_data()

10. 主控模块—main()

\section{系统设计}
\subsection{系统结构图}
系统采用主控函数统一调度、各功能模块独立实现的结构。main()启动时加载数据，进入菜单循环，根据用户选择分发到各功能函数，每次变更后自动保存数据到文件。

核心调用关系：main()→load\_data()→菜单循环(switch-case)→各功能函数→save\_data()

\subsubsection{数据结构设计}
\#define MAX\_PACKAGES 100

\#define MAX\_TRACKING  20

\#define MAX\_NAME      20

\#define MAX\_PHONE     15

\#define MAX\_COMPANY   20

\#define MAX\_SHELF     10

\#define MAX\_DATE      12

\#define MAX\_STATUS    8

\#define FILENAME      "packages.txt"

typedef struct \{

char tracking\_no[MAX\_TRACKING];

char receiver[MAX\_NAME];

char phone[MAX\_PHONE];

char company[MAX\_COMPANY];

char shelf[MAX\_SHELF];

char date[MAX\_DATE];

char status[MAX\_STATUS];

\} Package;

Package packages[MAX\_PACKAGES];

int package\_count = 0;

\subsection{模块要点}
核心算法说明：

1. 线性查找(find\_by\_tracking)：遍历数组逐项比较单号，O(n)

2. 快递录入：检查上限→唯一性→填充→追加→save\_data()

3. 删除快递：找索引→前移覆盖→count--→save\_data()

4. 数据统计：遍历数组，按status累加计数；按公司/货架去重

文件读写流程：

1. load\_data()：fopen→逐行fscanf解析→填充数组→fclose()

2. save\_data()：fopen→逐行fprintf写入→fclose()

3. 每次修改操作后自动save\_data()，退出前在case 0中调用

\section{功能演示与运行结果}
\subsection{功能测试用例}
\subsection{程序运行截图}
以下为系统各功能模块的运行效果说明。实际截图将在程序编译运行后补充插入。

\begin{figure}[H]
  \centering
  \fbox{\begin{minipage}{0.8\textwidth}
    \centering\vspace{3cm}
    \small[图1: 主菜单界面界面 -- 图片尚未插入]
    \vspace{3cm}
  \end{minipage}}
  \caption{图1: 主菜单界面界面}
\end{figure}


\begin{figure}[H]
  \centering
  \fbox{\begin{minipage}{0.8\textwidth}
    \centering\vspace{3cm}
    \small[图2: 快递录入界面 -- 图片尚未插入]
    \vspace{3cm}
  \end{minipage}}
  \caption{图2: 快递录入界面}
\end{figure}


\begin{figure}[H]
  \centering
  \fbox{\begin{minipage}{0.8\textwidth}
    \centering\vspace{3cm}
    \small[图3: 快递查询界面 -- 图片尚未插入]
    \vspace{3cm}
  \end{minipage}}
  \caption{图3: 快递查询界面}
\end{figure}


\begin{figure}[H]
  \centering
  \fbox{\begin{minipage}{0.8\textwidth}
    \centering\vspace{3cm}
    \small[图4: 取件管理界面 -- 图片尚未插入]
    \vspace{3cm}
  \end{minipage}}
  \caption{图4: 取件管理界面}
\end{figure}


\begin{figure}[H]
  \centering
  \fbox{\begin{minipage}{0.8\textwidth}
    \centering\vspace{3cm}
    \small[图5: 退件处理界面 -- 图片尚未插入]
    \vspace{3cm}
  \end{minipage}}
  \caption{图5: 退件处理界面}
\end{figure}


\begin{figure}[H]
  \centering
  \fbox{\begin{minipage}{0.8\textwidth}
    \centering\vspace{3cm}
    \small[图6: 信息修改界面 -- 图片尚未插入]
    \vspace{3cm}
  \end{minipage}}
  \caption{图6: 信息修改界面}
\end{figure}


\begin{figure}[H]
  \centering
  \fbox{\begin{minipage}{0.8\textwidth}
    \centering\vspace{3cm}
    \small[图7: 数据统计界面 -- 图片尚未插入]
    \vspace{3cm}
  \end{minipage}}
  \caption{图7: 数据统计界面}
\end{figure}


\begin{figure}[H]
  \centering
  \fbox{\begin{minipage}{0.8\textwidth}
    \centering\vspace{3cm}
    \small[图8: 浏览全部界面 -- 图片尚未插入]
    \vspace{3cm}
  \end{minipage}}
  \caption{图8: 浏览全部界面}
\end{figure}


\subsection{边界测试与健壮性测试}
\section{课程设计总结}
\subsection{收获与体会}
1. 结构化编程能力提升：将完整业务系统拆分为十余个功能函数，主控函数统一调度，深刻体会了“高内聚、低耦合”的模块化设计思想

2. 结构体实际应用：掌握struct定义、结构体数组组织方式，理解用结构体封装现实世界实体的建模方法

3. 文件I/O深入理解：熟练掌握fopen/fscanf/fprintf/fclose的使用，理解文件作为数据持久化存储的基本原理

4. 字符串处理技巧：学会strcmp(精确比较)、strstr(子串匹配)、strcpy(复制)等函数的应用场景和注意事项

5. 完整项目流程经验：从需求分析→模块设计→编码实现→测试验证，体验软件开发全生命周期

\subsection{遇到的问题与解决}
\subsection{优化或扩展的方向}
1. 动态内存管理：改为链表或动态数组(malloc/realloc)

2. 取件码功能：入库生成随机取件码

3. 排序功能：按时间/公司/货架排序

4. 逾期提醒：计算存放天数，超阈值标记

5. 图形化界面：EasyX或Qt

6. 多用户管理：操作员登录+操作日志

\section{参考文献}
[1] 谭浩强. C程序设计(第五版)[M]. 北京:清华大学出版社, 2017.

[2] Brian W.Kernighan,Dennis M.Ritchie. The C Programming Language(2nd Edition)[M]. Prentice Hall, 1988.

[3] 明日科技. C语言从入门到精通(第4版)[M]. 北京:清华大学出版社, 2019.

[4] ISO/IEC 9899:1999. Programming Languages—C[S]. ISO, 1999.

\section{附录：源程序清单}
\subsection{源程序文件名}
\subsection{编译与运行说明}
1. 编译环境:Windows 10/11+Dev-C++ 5.11(或GCC 5.0+)

2. 编译命令(GCC):gcc package\_system.c -o package\_system.exe -fexec-charset=GBK

3. 若使用Dev-C++则直接在IDE中打开package\_system.c编译运行即可

4. 运行:双击package\_system.exe或在命令行中执行

\subsection{完整源代码}